\section{Computations}\label{sec:computation}

\subsection{Arithmetic}\label{ssec:arith}

All inequalities that enter the proofs are verified in \emph{ball arithmetic}: every real number is represented by a midpoint and a
radius, and every operation returns a ball that contains the exact result. We use the Arb library \cite{Johansson2017Arb}, which is
part of FLINT \cite{flint} (version~3.6.0), through its Python interface \cite{pythonflint} (version~0.9.0, with CPython~3.14.2). The
endpoints of all boxes, the base points and weights of the supports and the integer data are exact rational numbers; in the lattice
case of Section~\ref{sec:bravais} some of them are binary floating-point numbers that are used as exact dyadic data. The kernels are
evaluated from
\[
  \Ga(\tfrac12,t)=\sqrt{\pi/t}\,\erfc\sqrt t,\qquad \Ga(1,t)=e^{-t}/t,\qquad \Ga(a+1,t)=\bigl(e^{-t}+a\,\Ga(a,t)\bigr)/t .
\]
An inequality is accepted only if it holds for the whole ball, and a computation that produces a non-finite ball is rejected.
Floating-point arithmetic occurs only in the programs that \emph{search} for certificates (to propose base points, weights,
rotations, the trial values $t$ of Section~\ref{ssec:taylor}, the coordinate to bisect and the order in which rules are tried)
and, in the verifying program of the cover, to order the triples of Lemma~\ref{lem:energy}(iii$'$); no inequality is accepted on
the basis of a floating-point computation.

The working precision is $128$ bits for the whole cover of Section~\ref{sec:cover}, $160$ bits in the local charts of
Section~\ref{sec:local}, $96$ bits in Section~\ref{sec:bravais} and $384$ bits for the constants. The precision affects
only whether an inequality can be accepted, never the validity of an accepted one. The smallest accepted margins are
about $10^{-12}$ in the cover and $10^{-10}$ in the lattice case. A program that checks the certificate must also check that its
parameters lie in the ranges that the lemmas require, for instance a positive splitting parameter $\alpha$ and supports built on
boxes with $r_1,r_2>0$: outside these ranges the bounds of Section~\ref{sec:cover} need not hold.

The proofs therefore depend on the following being correct: the ball operations of Arb used here (the arithmetic operations,
integer powers and roots, the exponential and the logarithm, the sine and cosine and their hyperbolic counterparts, the
complementary error function, the integer part, the constant $\pi$, and comparisons); exact integer and rational arithmetic in
Python; the programs that implement the rules of Sections~\ref{sec:bravais} and~\ref{sec:local}; and the program that verifies
the certificate of the cover (Section~\ref{ssec:certs}). The measures taken to limit this dependence are described in
Sections~\ref{ssec:checks} and~\ref{ssec:lean}.

\subsection{The constants}\label{ssec:constcert}

The energies of Proposition~\ref{prop:constants} are computed from \eqref{eq:ewald} for the primitive cells of the two cubic lattices,
with exact rational Gram matrices, for two different values of the splitting parameter ($\alpha=0.9$ and $\alpha=1.4$ in units of the
inverse Wigner--Seitz radius). The sums are extended until the exponents of the Gaussian factors exceed $160$, and the remainders are
bounded by the following elementary estimate: if the points of $P\subset\R^3$ have mutual distances at least $2\rho$, then the number
of points of $P$ in the ball of radius $r$ is at most $(1+r/\rho)^3$, and summation by parts bounds $\sum_{p\in P,|p|>r}\psi(|p|)$ for
decreasing $\psi$. At $384$ bits of precision the two computations give overlapping enclosures of width below $10^{-45}$, with
intersection $[\eB^-,\eB^+]$ and $[\eF^-,\eF^+]$, where
\begin{align*}
  \eB^-&=-0.89592925568176528916624546293829460538061046722070, & \eB^+&=\eB^-+4\times10^{-50},\\
  \eF^-&=-0.89587361519510845989748779611028534450181510608894, & \eF^+&=\eF^-+4\times10^{-50}.
\end{align*}
These are the numbers used in Section~\ref{sec:cover}, and $\gap^+=\eF^+-\eB^-$. They agree with the twelve-decimal values in the
literature \cite{BaikoPotekhinYakovlev2001} and are contained in the enclosures obtained in \cite{MuratovNovaga2026} by a different
method and in a different arithmetic: for $\eB$ in \cite[Lemma~2]{MuratovNovaga2026} and for $\eF$ in the remark after its proof,
where they are stated for the dual lattices (see \cite[(7), (17), (18)]{MuratovNovaga2026} for the conversion).

\subsection{The certificates}\label{ssec:certs}

\emph{The lattice case.} The program for Theorem~\ref{thm:bravais} is deterministic. It examines the $1.78$ million boxes of
Table~\ref{tab:bravais} in about three hours on four processor cores and records the number of boxes of each type and the smallest
margins; the boxes themselves are not stored, and the verification consists in running the program.

\emph{The local charts.} For each of the eight ellipsoids of Table~\ref{tab:atlas} the certificate records the enclosures of the
Hessian at the centre, the certified shift $\lambda$, the coefficients of the cubic polynomial, the bounds $B_3$, $B_4$ with their
tails, and the geometric conditions of Section~\ref{sec:local}; it is recomputed in three minutes.

\emph{The cover.} The certificate consists of thirteen files in a documented text format, $187$~MB in compressed
form. It contains, for every box received by a stage, its subdivision and, for every part that is not held, the rule
and the exact data that determine its parameters (Section~\ref{ssec:stages}). The certificate is verified as it stands, by a
single program, in two passes.
\begin{itemize}
\item[(a)] \emph{Completeness}, in exact rational arithmetic: every bisection has both halves; every other subdivision
is tested to be a cover by an exact recursive procedure; the boxes replaced by a file are exactly the boxes held by the
preceding one; the reflection step is carried out; every containment of type (c) is tested; the first file starts from
$\Kbox$ and the last one holds nothing. No volume is computed.
\item[(b)] \emph{Inequalities}, in ball arithmetic: for every part that carries a rule and for every box that a file lists
with a rule (type (c) of Section~\ref{ssec:stages}), the hypothesis of the rule is recomputed on that box from its parameters.
\end{itemize}
The recorded run of this program over the whole certificate, at $128$ bits, took $18.6$ processor-hours ($3.3$ hours on $6$
processes) on the workstation described below and accepted all $17\,092\,214$ parts and boxes with a rule
(Table~\ref{tab:stages}).

The recorded runs of the programs and of the verifier, which accompany the paper (Section~\ref{ssec:avail}), were made on one
workstation with $32$ logical processors and $62$~GiB of memory.

\subsection{Checks}\label{ssec:checks}

The program that verifies the cover implements the statements of Sections~\ref{sec:setting} to~\ref{sec:cover} and of
Appendices~\ref{app:tails} and~\ref{app:details} that the certificate needs. It consists of fourteen Python modules (about
$4\,300$ lines) and uses only the standard library and Arb; it imports no other code. What it takes
as given is Theorem~\ref{thm:bravais}, the enclosures of $\eB$ and $\eF$, Montgomery's theorem and the lemmas of this paper;
it checks every inequality of the certificate and the coverage of $\Kbox$ as described in Section~\ref{ssec:certs}. The program is accompanied by negative
controls: every rule is applied, through the same entry point, to boxes around a body-centred and a face-centred cubic
point, where (N), (E) and (D) are false, with parameters throughout and outside their ranges; parts of the certificate are
changed in one parameter or moved outside the hypotheses of their rule; and boxes are removed from the certificate, covers
and containments falsified, and the reflection step misused. All controls are rejected, and in every control that changes
a part of the certificate the unchanged part is accepted.

The constants agree with the independent computation of \cite{MuratovNovaga2026} (Section~\ref{ssec:constcert}). The verifier
of the cover and the programs for the lattice case, the constants and the local charts use the same ball arithmetic library,
so these checks reduce the risk of programming errors but are not independent of the trust assumptions of
Section~\ref{ssec:arith}.

\subsection{Formal verification}\label{ssec:lean}

A formalization in Lean~4 (version 4.35.0) with Mathlib accompanies the paper; the part described here consists of $405$
modules. No axioms are used beyond the three standard axioms of Lean's logic with classical choice (propositional
extensionality, quotient soundness and choice), and no proof relies on compiled evaluation: all finite computations are
evaluated by the kernel. The following is proved in it without hypotheses.
\begin{itemize}
\item \emph{Definitions and invariances.} The periodic Coulomb kernel and its Fourier coefficients; the equality of the energy
\eqref{eq:edef} with Ewald's expression for every splitting parameter and with the positive series \eqref{eq:H1}, \eqref{eq:H2}; its
independence of the period lattice and its invariance under isometries; the definition of the defect, its invariances, and its
equality at all points of a two-site configuration.
\item \emph{Constants.} Enclosures of $\eB$, $\eF$ and $\gap$ to fifty digits, consistent with Section~\ref{ssec:constcert}.
\item \emph{The lattice case.} For every lattice of density $3/(4\pi)$, inequality \eqref{eq:bravais} holds, and equality implies
that the lattice is BCC or FCC; moreover $D=0$ on BCC and $D=1$ on FCC. This is Theorem~\ref{thm:bravais}. The formal proof does not
replay the certificate of Section~\ref{sec:bravais}; it uses a second, smaller certificate designed for evaluation in the kernel,
with tangent-line bounds in fixed-point rational arithmetic: a cover of the reduced domain by $5\,671$ boxes and covers of two
strain charts by about $8\,700$ boxes. Theorem~\ref{thm:bravais} thus has two computer-assisted proofs with different certificates
and different arithmetic, one of them checked by a proof assistant.
\item \emph{Analytic part of the two-site theorem.} The compact reduction (Proposition~\ref{prop:compact}); the existence of a
registration with the inequality \eqref{eq:registered} for $D<1$ and $e\le\eF$ (but not the uniqueness statements of
Theorem~\ref{thm:registration}); upper bounds for the defect in terms of a given registration, as in Lemmas~\ref{lem:Dupper}
and~\ref{lem:Dbox}, and bounds for its growth away from an anchor, as in Lemma~\ref{lem:growth}; and Montgomery's inequality for the
planar theta function \cite{Montgomery1988}. The local inequality of Section~\ref{sec:local} is not formalized.
\item \emph{Checking functions.} For most kinds of sufficient conditions used in the cover, a function with a proof of soundness:
if the function accepts a box with its data, then the box is discharged. These functions were written for the formalization; several
of them use bounds different from those of Section~\ref{sec:cover}.
\item \emph{Reduction of the theorem.} A theorem that derives, from the hypothesis that the checking functions accept the data of
the cover, the conjunction of Theorems~\ref{thm:min}, \ref{thm:main} and~\ref{thm:basin} (the last with $\|A-U\|_{\mathrm F}$ for some
rotation $U$ in place of $\dist(A,\SO(3))$, which it bounds from above, so that this form implies Theorem~\ref{thm:basin});
Corollary~\ref{cor:fcc} follows from it in a few lines.
\end{itemize}
What is \emph{not} done is to establish this hypothesis, that is, to evaluate checking functions on the certificate of the cover
inside the kernel; the modules evaluate checking functions in the kernel only on samples of boxes of several stages. For
some kinds of anchor and deficit comparisons no formal checking function is written, and some of the existing ones are not
strong enough to accept the data of the certificate, so that the missing step needs more than computing time. Consequently Theorem~\ref{thm:cover}, and with it Theorems~\ref{thm:min}--\ref{thm:basin} beyond the
lattice case, are not formally verified: they rest on the ball-arithmetic verification described above.

The $405$ modules of the package are compiled in a single clean build, each by its own call of Lake in an order in which every
module comes after its imports, from one set of sources, identified by their SHA-256 in the records, and with Mathlib at one
revision; the recorded build takes
$3.8$ hours of elapsed time and $12.6$ processor-hours, and Lean's kernel checks every theorem of the modules in it. The
compiled modules are then replayed with the checking program distributed with Lean (\texttt{leanchecker}), which reads each
compiled module and checks every declaration of it with the kernel; it accepts all of them. The $130$ theorems that represent
the items above are exported with all the declarations they depend on, Mathlib included, and checked from the export by
Lean's kernel and by Lean4Lean, a re-implementation of the kernel in Lean, both of which accept all of them; the $54$ of them
that depend on no kernel-evaluated certificate are checked in addition by nanoda and by con-leche and its Rust counterpart
con-ron. The export declares exactly the three axioms named above, and an audit of all $47\,462$ constants of the $405$
modules finds no other. The Lean sources and the records of the build and of these checks are part of the package described
in Section~\ref{ssec:avail}.

\subsection{Availability}\label{ssec:avail}

The certificates, the verifying program of the cover, the programs for the lattice case, the constants and the local charts,
and the Lean sources described in Section~\ref{ssec:lean} (Lean~4.35.0, Mathlib at revision \texttt{56a00585}), with
instructions for reproducing every computation, form a package that accompanies this paper. It is available at
\url{https://github.com/ChaseCrawff/bcc-sharp}.
